\section[PHT3D Exercise 3: Precipitation/dissolution fronts in acid mine drainage]{PHT3D Exercise 3: Precipitation/dissolution fronts in acid mine drainage}

%\subsection{Introduction}

This exercise is another simulation problem that involves mineral dissolution and precipitation reactions as the principle geochemical reactions. Compared to the previous 
exercise it is somewhat more complicated as it contains significantly more 
aqueous components and mineral reactions.
It was first presented by \cite{walt94a} and later also used as a benchmark problem by \cite{gmt3d}. In contrast to the previous exercise this case includes now also redox reactions. 
The simulations demonstrate the typical hydrogeochemical changes that occur 
when acidic mine tailings leach into an anaerobic carbonate aquifer. 
Aqueous complexation and dissolution/precipitation are all considered as equilibrium reactions. 
If the original reaction network defined by \cite{walt94a} is used, the simulation 
includes 17 aqueous components, 15 of which are transported, 54 aqueous species and six minerals. 

\subsection[Spatial discretisation and flow problem]{Spatial discretisation and flow problem}

The problem is set up as a simple two-dimensional flow and transport problem with dimensions and properties as defined in \ref{ex3_flow}. To construct the MODFLOW model that underlies the transport problem proceed as follows:

%\begin{figure}[b]
%\centering
%\includegraphics[width=80mm]{walter.eps}
%\caption[Initial definition of the model grid for PHT3D Exercise 2]{Initial definition of the model grid for PHT3D %Exercise 2}
%\label{fig:mesh_walter}
%\end{figure}


\begin{itemize}

\item In ORTI, create a new model and save it again in a new, separate directory.

\item Under {\color{blue}Parameters} $\rightarrow$ {\color{blue}Model} select the {\color{blue}Model} button and specify the attributes of the new model. Select the values {\color{blue}Xsection} for {\color{blue}Dimension}, {\color{blue}Confined} for {\color{blue}Type} 
and {\color{blue}MODFLOW Family} for {\color{blue}Group} 

\item Under {\color{blue}Parameters} $\rightarrow$ {\color{blue}Model} select the {\color{blue}Grid} button and specify the details of the model domain and the spatial discretisation in the dialog that opens up. Define a vertical transect that is 100m long and 10m thick. The model domain
may be discretised into grid cells of 4m length and 1m height. Note, that the 
model discretisation can be easily changed again later. 

\item Select {\color{blue}Parameters} and select the {\color{blue}Time} button.
Set the {\color{blue}Total Simulation Time} to 2000 days and select a time {\color{blue}Step size}
of 20 days. 

\item Go to {\color{blue}Spatial Attributes} $\rightarrow$ {\color{blue}MODFLOW} $\rightarrow$ {\color{blue}LPF} $\rightarrow$ {\color{blue}lpf.8} to specify the hydraulic conductivity to 1 $m/day$

\item Go to {\color{blue}Spatial Attributes} $\rightarrow$ {\color{blue}MODFLOW} $\rightarrow$ {\color{blue}LPF} $\rightarrow$ {\color{blue}lpf.9} to specify the vertical hydraulic conductivity to 1 $m/day$

\item Go to {\color{blue}Spatial Attributes} $\rightarrow$ {\color{blue}MT3DMS} $\rightarrow$ {\color{blue}BTN} $\rightarrow$ {\color{blue}btn.11} and specify the effective porosity to be 0.35.


\item Implement hydraulic boundary conditions by defining the hydraulic heads at the upstream and the downstream boundary of the model. First, define the type of boundary condition by going to {\color{blue}Spatial Attributes} $\rightarrow$ {\color{blue}MODFLOW} $\rightarrow$ {\color{blue}BAS6} $\rightarrow$ {\color{blue}bas.3 Boundary Conditions} and setting the value to 1 (= active cell) for all cells. After that select {\color{blue}Zone} and use the {\color{blue}Zone Tools} to add one line at the inflow end of the model and one line at the effluent end. 
In both cases the lines need to extend over the whole depth of the aquifer. 
The value for both lines (=zones) needs to be set to -1. Use, for example,
$bc\_up$ and $bc\_down$ as {\color{blue}Zone Names}. You will see the selected names displayed near 
the location of the lines.
  
This will define the type of boundary condition to a fixed head boundary 
at the inflow and outflow model boundaries.
The hydraulic head at these grid cells will therefore remain unchanged
over the whole simulation period. It will remain at the same values that we will now define as initial heads.  

\item Set initial hydraulic heads by going to {\color{blue}Spatial Attributes} $\rightarrow$ {\color{blue}MODFLOW} $\rightarrow$ {\color{blue}BAS6} $\rightarrow$ {\color{blue}bas.6 Initial Heads} and set {\color{blue}one\_value} 
to 10, corresponding to a head of 10m. Then press {\color{blue}'OK'} to confirm the value. In addition
define a zone (line) at the upstream end of the model domain (where you also defined that the type of 
boundary condition is -1) and allocate a hydraulic head of 12m over the entire depth of the aquifer. 

% Maybe lets not use recharge here to save time
%\item Add a value for recharge by selecting Models $\rightarrow$ MODFLOW $\rightarrow$ Flow 
% Packages $\rightarrow$ Recharge $\rightarrow$ Reset Matrix and change the value for the Recharge % flux to 0.001 m/d (365 mm/year) and click ok. 

\end{itemize}

\begin{table}[t]
\caption{Flow and transport parameters for the simulation.}
\begin{center}
\begin{tabular}{lc} 
\toprule
Flow simulation                     & steady state \\
Total simulation period\ $ (days)      $         &  2000            \\
Model \  length\ $ (m)      $         &  100            \\
Model \  thickness\ $(m)$                        &  10            \\
Grid \ spacing\ $\Delta x $ $(m)$                &    4            \\
Grid \ spacing\ $\Delta z $ $(m)$                &    1            \\
Porosity\     $            $         &  0.35           \\
Horizontal hydraulic conductivity\ $(m/day)$               &  1         \\ 
Vertical   hydraulic conductivity\ $(m/day)$               &  1         \\ 
Piezometric head upstream boundary\ $(m)          $         &  12           \\ 
Piezometric head downstream boundary\ $(m)        $         &  10           \\ 
Longitudinal dispersivity\ $(m)         $         &  2.5          \\ 
Transverse  vertical dispersivity\ $(m)         $         &  0.025          \\ 
\bottomrule
\end{tabular}
\label{ex3_flow}
\end{center}
\end{table}

\subsection{Running MODFLOW}

You have now completed to set up the flow model and you are ready to run MODFLOW.

\begin{itemize}

\item Under {\color{blue}Parameters} $\rightarrow$ {\color{blue}Flow} use the {\color{blue}Write Button} to write all the input files required to run MODFLOW

\item Under {\color{blue}Parameters} $\rightarrow$ {\color{blue}Flow} use the 
{\color{blue}Red Arrow} button to start the MODFLOW simulation

\item Check the simulated head contours and assess whether the results are plausible.

% Olivier, maye we can practise the use of some more iPHT3D output features 

\end{itemize}

% Here we had a non-reactive model step in the pmwin version of the tutorial

\subsection[Data input for solute transport properties]{Data input for solute transport properties}

For the simulation of the solute transport of the mobile components the model parameters that control advective and dispersive transport need to be entered. 

\begin{itemize}

\item Go to {\color{blue}Parameters} $\rightarrow$ {\color{blue}Transport} and select the {\color{blue}Parameter} button

\item In the dialog box that subsequently appears, select {\color{blue}ADV} and select under {\color{blue}adv.1 Solver Parameters} the {\color{blue}3rd-order TVD Scheme (ULTIMATE)} as the Solution Scheme. Leave the default parameters and click {\color{blue}ok}.

\item Select {\color{blue}Spatial Attributes} $\rightarrow$ {\color{blue}MT3DMS} $\rightarrow$ {\color{blue}DSP}, enter a longitudinal dispersivity value of 2.5 m. and click {\color{blue}OK}

\item Select {\color{blue}Spatial Attributes} $\rightarrow$ {\color{blue}MT3DMS} $\rightarrow$ {\color{blue}DSP}, and specify that the ratio between transverse vertical dispersivity and the longitudinal dispersivity is 0.01.
This will define that the actual value is 0.025 m (= 25mm). 


\end{itemize}

\subsection[Data input for reactive transport]{Data input for reactive transport}

In this step we define the reaction network and the various water and mineral compositions that play a role in this problem. Both the initial and the inflow concentrations of of the aqueous components that will be used in the simulation are shown in Table \ref{ex3_caqu}. Also the initial concentrations of the minerals are
given in Table \ref{ex3_cmin}. 

\begin{table}
\caption{Aqueous concentrations used in the Walter et al. Example.}
\begin{center}
\begin{tabular}{l c c}
\hline
Aqueous  component        & $C_{init}$ & $C_{tailing}$      \\
                  & $(mol\ l^{-1}_w)$   & $(mol\ l^{-1}_w)$      \\ \hline
pH              &   6.96    &  3.99           \\
pe              &   1.67    &  7.69           \\
C(4)            &  3.94 $\times$ $10^{-3}$ & 4.92 $\times$ $10^{-4}$   \\
S(6)            &  7.48 $\times$ $10^{-3}$ & 5.00 $\times$ $10^{-2}$   \\
S(-2)           &     -             &     -              \\
Fe(2)           &  5.39 $\times$ $10^{-5}$ & 3.06 $\times$ $10^{-2}$   \\
Fe(3)           &  2.32 $\times$ $10^{-8}$ & 1.99 $\times$ $10^{-7}$   \\
Mn(2)           &  4.73 $\times$ $10^{-5}$ & 9.83 $\times$ $10^{-6}$   \\
Ca              &  6.92 $\times$ $10^{-3}$ & 1.08 $\times$ $10^{-2}$   \\
Mg              &  1.96 $\times$ $10^{-3}$ & 9.69 $\times$ $10^{-4}$   \\
Na              &  1.30 $\times$ $10^{-3}$ & 1.39 $\times$ $10^{-3}$   \\
K               &  6.65 $\times$ $10^{-5}$ & 7.93 $\times$ $10^{-4}$   \\
Cl              &  1.03 $\times$ $10^{-3}$ & 1.19 $\times$ $10^{-4}$   \\
Al              &  1.27 $\times$ $10^{-7}$ & 4.30 $\times$ $10^{-3}$   \\
Si              &  1.94 $\times$ $10^{-3}$ & 2.08 $\times$ $10^{-3}$   \\ \hline
\end{tabular}
\label{ex3_caqu}
\end{center}
\end{table}

\begin{table}
\caption{Mineral concentrations Walter et al. Example.}
\begin{center}
\begin{tabular}{ l c }
\hline
Mineral        & $C_{init}$  \\
                  & $(mol\ l^{-1}_{v})$  \\ \hline
Calcite \   $(CaCO_3)  $  & 1.95 $\times$ $10^{-2}$  \\
Siderite  \ $(FeCO_3)  $  & 4.22 $\times$ $10^{-3}$  \\
Gibbsite \ $(Al(OH)_3) $  & 2.51 $\times$ $10^{-3}$  \\
amorphous \ $Fe(OH)_3  $  & 1.86 $\times$ $10^{-3}$  \\
Gypsum \ $(CaSO_4.2H_2O) $ & 0.0                     \\
amorphous  \ $SiO_2    $  & 4.07 $\times$ $10^{-1}$  \\ \hline
\end{tabular}
\label{ex3_cmin}
\end{center}
\end{table}

\begin{itemize}

\item Copy the database file that is going to be used for the PHT3D simulation
into the folder that contains all files for this exercise and name the file {\color{red} \textbf{pht3d\_datab.dat}}.
You can use the same database as in PHT3D Exercise 1, i.e., the standard PHREEQC database.

\item Go to {\color{blue}Parameters} $\rightarrow$ {\color{blue}Chemistry} and select the
{\color{blue}Import database} button. ORTI will then read and interpret 
the PHT3D database file.

\item In this next step we define all initial concentrations of all aqueous components and minerals and 
also the concentrations at the upstream model boundary (which will be allocated to the inflowing groundwater).

The composition of the ambient water (initial water composition at the start of the simulation) 
needs to be entered under the {\color{blue}Solutions} Tab in the {\color{blue}Background} column.
Activate all relevant species by ticking the appropriate boxes for 
$pH$, $pe$, $C(4)$, $S(6)$, $S(-2)$, $Fe(2)$, $Fe(3)$, $Mn(2)$, $Ca$, $Mg$, $Na$, $Cl$, $Al$ and $Si$
and enter the concentrations. The set of minerals that are considered in this simulation
must be activated and the corresponding mineral composition must be entered 
under the {\color{blue}Phase} Tab.

In addition to the background water composition we will also need to define the composition of the 
acidic tailings water that enters the model domain in the upper section of the upstream boundary. 
This water composition will be defined in the {\color{blue}Solution 1} column 
under the {\color{blue}Solutions} Tab.
For the lower section of the inflow boundary it is assumed that the inflowing water
is the same as the background water composition.   

\item The allocation of the appropriate water compositions as model boundary condition can be
done by selecting {\color{blue}Spatial Attributes} $\rightarrow$ {\color{blue}MT3DMS} $\rightarrow$ {\color{blue}BTN} $\rightarrow$ {\color{blue}btn.12 Boundary Condition} and defining a zone (line)
for the lower section of the inflow (upstream) boundary (between y = 0m and y = 7m)
for which a value of -1 is allocated. This defines that the concentrations at these grid cells 
will not change during the entire simulation.   
The concentrations that ORTI will allocate to the model's grid cells are the ones defined 
as initial concentrations. This means that the water composition that was defined as background
water composition will also be used as inflow water composition. 

To additionally define that the acidic leachate composition is used in the upper section we need to 
select {\color{blue}Spatial Attributes} $\rightarrow$ {\color{blue}PHT3D} $\rightarrow$ {\color{blue}PH} $\rightarrow$ {\color{blue}ph.4 Source \/ Sink} and define a zone (line) for the uppermost 3 m (y= 7m to y = 10m)
of the boundary (the depth zone at which the acidic leachate enters) and allocate 
a {\color{blue}Zone Value} of 1. This induces that the concentrations that were previously defined 
for $Solution\ 1$ in the {\color{blue}Solutions} Tab will be applied to the water that enters 
the model domain in the upper section of the inflow boundary. 

\end{itemize}

\subsection[Running PHT3D]{Running PHT3D}

To run PHT3D, proceed as earlier in the exercise when executing MODFLOW:

\begin{itemize}

\item Go to {\color{blue}Parameters} $\rightarrow$ {\color{blue} Chemistry} and use 
the {\color{blue}Write Button} to write all the input files required to run PHT3D

\item Then go to {\color{blue}Parameters} $\rightarrow$ {\color{blue}Chemistry} and use the {\color{blue}Plume Button} to start PHT3D.

\end{itemize}

\subsection[Visualisation of results]{Visualisation of results}

After the end of the simulation visualise the results in various ways.

\begin{itemize}

\item Obtain a contour plot by navigating in the {\color{blue} Results} panel of ORTI.
Select a specific timestep and a species for which you want to see your results.
For example to see the pH contours after 1000 days select {\color{blue}Results} $\rightarrow$ {\color{blue}Aquifer} 
$\rightarrow$ {\color{blue}Tstep} and select {\color{blue}1000} and then go to   
{\color{blue}Results} $\rightarrow$ {\color{blue}Chemistry} 
$\rightarrow$ {\color{blue}Species} and select {\color{blue}pH}.  
You can easily move forward and backward in time by clicking {\color{blue}Tstep} once more and then
using your keyboard arrows to move forward or backward.

\item Plot a breakthrough curve at a selected location to visualise how concentrations change over time.
First, define the location for which you want to plot breakthrough curves by
going to {\color{blue}Spatial Attributes} $\rightarrow$ {\color{blue}Observation} $\rightarrow$ {\color{blue}OBS} $\rightarrow$ {\color{blue}obs.1} $\rightarrow$ {\color{blue}Zone}
and use your mouse to position the cursor and clicking at the location
where you wan to create an observation point. In the dialog that comes up name the zone $BTC1$. 
No specific input is required for the {\color{blue}Zone Value} field.
The coordinates of the location selected by the mouse click can be modified, if required.  
 
\end{itemize}




\subsection[Discussion of simulation results]{Discussion of simulation results}

\par Only a few aspects of the geochemical evolution
observed in the model simulation are described
in the following. For a detailed description see \cite{walt94a}.
The acidic inflow solution is initially buffered by 
calcite ($CaCO_3$), maintaining the $p$H
at an almost neutral level (6.5 - 7). In this zone, gypsum
($CaSO_4.2H_2O$) is formed from calcium released during calcite
dissolution and sulphate ($S(6)$) from the inflow solution
($C_{in,S(6)} > C_{backgr,S(6)}$). Similarly, siderite ($FeCO_3$)
can form from $C(4)$ ($CO_3$) released during dissolution and the
$Fe(2)$-rich inflow solution. At locations where calcite is
completely dissolved, siderite becomes the buffering mineral,
dissolving until it is also entirely removed. Finally, gibbsite
($Al(OH)_3$) precipitation is the buffer mechanism at locations
where both calcite and siderite are completely removed. The three
different buffering mechanisms lead to three distinct levels of
$p$H. Also three distinct
$p$e levels evolve with fronts from more reduced to more oxidised
water. The fronts have the same positions as those in the $p$H
plot. As the $p$e of the solution is completely controlled by the
redox couple $Fe(2)/Fe(3)$, the $p$e changes reflect the changes
in their concentration ratio. Both $Fe(2)$ and $Fe(3)$
concentrations are strongly controlled by the $p$H of the
solution. Generally, their solubility increases with a decreasing
$p$H. However, as the increase of their solubilities is not
proportional, the $Fe(2)/Fe(3)$ ratio can change and so does the
$p$e.

\subsection[Model variants]{Model variants}

\begin{itemize}

\item To see the effect of calcite buffering rerun the model with a much 
smaller initial calcite concentration and compare how far the 
low pH zone and dissolved aluminium are migrating in this case.

\item Test the effect of calcite-filled permeable reactive barrier (PRB).
To do that add a small zone (between x = 40m and x = 50 m and y = 0m and y = 10m) 
where a high concentration of calcite is present but no other minerals. 
After rerunning the model inspect the evolution of the low pH zone and of the aluminium migration.

\end{itemize}

